\chapter{Three Generations: Why Not Four?}
\label{ch:generations}

\begin{quote}
\textit{This chapter addresses one of the deepest puzzles in particle physics:
why does the fermion spectrum consist of exactly three generations?
EDC proposes a geometric answer rooted in $\mathbb{Z}_6$ crystallography.}
\end{quote}

% ════════════════════════════════════════════════════════════════════════════════
\section{The Generation Puzzle}
\label{sec:gen_puzzle}
% ════════════════════════════════════════════════════════════════════════════════

\begin{tcolorbox}[edcGuardrail, title=\textbf{Epistemic Status: HIGH RISK}]
\textbf{What we have:}
\begin{itemize}[nosep]
    \item An identification linking $\mathbb{Z}_3 \subset \mathbb{Z}_6$ to generation count \tagI{}
    \item Mode indices $n = 0, 1, 2$ mapped to $(e, \mu, \tau)$ with $<1\%$ mass error \tagI{}
    \item Multiple candidate mechanisms for truncation (none fully derived)
\end{itemize}

\textbf{What remains open:}
\begin{itemize}[nosep]
    \item A rigorous derivation of \emph{why} exactly three modes survive
    \item The dynamical mechanism forbidding $n \geq 3$ modes
\end{itemize}

\textbf{Honest assessment:} This chapter documents an \textbf{identification}, not a
\textbf{derivation}. The structural pattern is compelling; the proof is incomplete.
\end{tcolorbox}

In the Standard Model, the existence of three fermion generations is an
\emph{empirical input}. The gauge structure $SU(3)_C \times SU(2)_L \times U(1)_Y$
permits any number of generations; the value $N_{\text{gen}} = 3$ is simply
observed \tagBL{}.

EDC aims to derive $N_{\text{gen}} = 3$ from geometric structure:
\begin{equation}
    \mathbb{Z}_6 = \mathbb{Z}_2 \times \mathbb{Z}_3
    \quad\Longrightarrow\quad
    |\mathbb{Z}_3| = 3 \;\stackrel{?}{\longleftrightarrow}\; N_{\text{gen}} = 3
    \label{eq:z6_factor}
\end{equation}

\begin{tcolorbox}[colback=yellow!5, colframe=yellow!50!black, title=\textbf{Dictionary: What carries $\mathbb{Z}_3$ charge?}]
\textbf{Open question} \tagOpen{}: The physical object carrying $\mathbb{Z}_3$ index
could be:
\begin{itemize}[nosep]
\item Angular sector of localized wavefunction $\psi(\xi, \phi)$ with $\phi \in \{0, 2\pi/3, 4\pi/3\}$
\item Eigenmode index $n \in \{0,1,2\}$ of a BVP tower
\item Topological winding class in the brane's internal space
\item Superselection sector from boundary conditions
\end{itemize}
\textbf{Status}: This is an \textbf{identification of structure} \tagI{}, not a derivation
of generation count. The ``$=$'' in $|\mathbb{Z}_3| = N_{\text{gen}}$ is a matching,
not a proof.
\end{tcolorbox}

The challenge is to make this connection rigorous, not merely suggestive.

% ════════════════════════════════════════════════════════════════════════════════
\section{Physical Picture: Three Angular Channels}
\label{sec:physical_picture}
% ════════════════════════════════════════════════════════════════════════════════

\begin{mechanism}{Generation Structure from Angular Localization}
\textbf{Step 1: The brane has internal structure.}
The EDC membrane has finite thickness $\delta$ along the fifth dimension $\xi$,
and internal angular structure. This structure is the key to generation physics \tagP{}.

\textbf{Step 2: Flux tubes arrange hexagonally.}
Energy minimization forces flux tubes to arrange in a hexagonal lattice---the
densest 2D circle packing. This lattice has $\mathbb{Z}_6$ rotational symmetry \tagP{}.

\textit{Note}: The claim ``energy minimization $\to$ hexagonal'' is plausible but
not derived from the EDC action. Required: show that the relevant functional
(membrane + flux tube energy) is minimized by hexagonal packing.

\textbf{Step 3: $\mathbb{Z}_6$ factorizes into $\mathbb{Z}_2 \times \mathbb{Z}_3$.}
\begin{itemize}[nosep]
    \item $\mathbb{Z}_2$: Matter/antimatter distinction \tagP{}
    \item $\mathbb{Z}_3$: Labels three inequivalent angular sectors \tagP{}
\end{itemize}

\textbf{Step 4: Fermions localize in one of three channels.}
Each fermion ``chooses'' one of three angular sectors at positions
$\phi = 0, 2\pi/3, 4\pi/3$. These become the three generations:
\begin{itemize}[nosep]
    \item Sector 0: Electron family ($e$, $\nu_e$, $u$, $d$)
    \item Sector 1: Muon family ($\mu$, $\nu_\mu$, $c$, $s$)
    \item Sector 2: Tau family ($\tau$, $\nu_\tau$, $t$, $b$)
\end{itemize}

\textbf{Step 5: Overlap between sectors controls mixing.}
The overlap integral between wavefunctions in different angular sectors
determines mixing matrix elements. Large separation (deep wells) produces
small off-diagonal elements, explaining the near-diagonal CKM structure \tagI{}.

\textbf{Step 6: Why only three?} That's the open question (see \S\ref{sec:mechanisms}).
\end{mechanism}

% ════════════════════════════════════════════════════════════════════════════════
\section{Candidate Mechanisms}
\label{sec:mechanisms}
% ════════════════════════════════════════════════════════════════════════════════

\subsection{Mechanism A: $\mathbb{Z}_3$ from Hexagonal Symmetry}

The hexagonal lattice guarantees $\mathbb{Z}_6$ rotational invariance \tagDc{}.
The factorization $\mathbb{Z}_6 = \mathbb{Z}_2 \times \mathbb{Z}_3$ is pure
mathematics \tagM{}. The proposal:
\begin{itemize}[nosep]
    \item $\mathbb{Z}_2$: Matter/antimatter (C-parity) \tagP{}
    \item $\mathbb{Z}_3$: Generation index (three-fold rotational symmetry) \tagP{}
\end{itemize}

\begin{center}
\begin{tabular}{lcl}
\toprule
\textbf{Criterion} & \textbf{Status} & \textbf{Comment} \\
\midrule
$\mathbb{Z}_6$ from hexagons & GREEN & 2D packing + energy min \tagDc{} \\
$\mathbb{Z}_6 = \mathbb{Z}_2 \times \mathbb{Z}_3$ & GREEN & Pure group theory \tagM{} \\
$|\mathbb{Z}_3| = 3 \to N_{\text{gen}}$ & RED & Cardinality matching, not derivation \\
Explains mode truncation & RED & Doesn't explain why $n \geq 3$ forbidden \\
\bottomrule
\end{tabular}
\end{center}

\textbf{Verdict:} YELLOW---The link is \textbf{identified} \tagI{}, not \textbf{derived}.

\subsection{Mechanism B: KK Tower Truncation}

In Kaluza-Klein theories, compactification produces an infinite tower of modes.
The proposal: only $n = 0, 1, 2$ survive because higher modes are unstable.

If higher modes tunnel through a barrier, the lifetime scales as:
\begin{equation}
    \tau_n \propto \exp\left(+\frac{S_n}{\hbar}\right)
    \quad\text{where}\quad
    S_n = \int_0^{\xi_*} \sqrt{2m_{\text{eff}}(V(\xi) - E_n)} \, d\xi
    \label{eq:lifetime}
\end{equation}

\textbf{Definitions needed} \tagOpen{}:
\begin{itemize}[nosep]
\item $m_{\text{eff}}$: effective mass in $\xi$-direction (from 5D kinetic term reduction)
\item $V(\xi)$: effective potential from brane profile (not yet derived from EDC action)
\item $\xi_*$: classical turning point where $V(\xi_*) = E_n$
\item Domain: $\xi \in [0, \xi_*]$ (or $[0, R_\xi]$ for compact extra dimension)
\end{itemize}

For $n = 3$ to be unobservable, we need $\tau_3 < \tau_{\text{obs}}$ where the
\textbf{observability threshold} is:
\begin{equation}
\tau_{\text{obs}} \sim 10^{-23} \text{ s} \quad \text{(hadronization / strong scale)} \quad \tagBL{}
\end{equation}
This is the relevant scale, not the Planck time $t_P \sim 10^{-44}$ s.

\textbf{Verdict:} RED---Plausible mechanism but \textbf{entirely uncomputed}.
Requires deriving $V(\xi)$ and $m_{\text{eff}}$ from the EDC action (\OPR{12}).

% BACKFILL-GAP14: μ-window constraint BEGIN
\subsubsection{The $\mu$-Window Constraint}

For three generations to be observable while $n \geq 3$ is forbidden, the
\textbf{barrier parameter} $\mu$ must satisfy a specific range:

\begin{mechanism}{Barrier Parameter Constraint}
\textbf{Definition} \tagDer{}:
\begin{equation}
\mu \equiv \int_0^{\xi_*} \frac{m_{\text{eff}}(\xi')}{\Lambda} \, d\xi'
\label{eq:mu_barrier}
\end{equation}
where $\Lambda$ is the characteristic mass scale from the brane profile.

\textbf{Stability window:}
For mode $n$ to be observable, its tunneling action must satisfy:
\begin{equation}
S_n > \hbar \ln\left(\frac{\tau_{\text{obs}}}{t_0}\right)
\quad\Rightarrow\quad
\mu_n > \mathcal{O}(1)
\end{equation}

\textbf{The ``$\mu$-window''} \tagP{}:
\begin{itemize}[nosep]
\item $\mu_0, \mu_1, \mu_2 \gtrsim 5$: First three modes stable (observable)
\item $\mu_3 \lesssim 1$: Fourth mode unstable (decays before observation)
\end{itemize}

\textbf{What this requires:}
The barrier profile $V(\xi)$ must produce a ``cliff'' between $n = 2$ and $n = 3$
where the effective barrier drops sharply. This is the content of OPR-12.
\end{mechanism}

\textbf{Status:} The $\mu$-window is the \textbf{closure target} for Mechanism B.
If EDC dynamics can derive $V(\xi)$ such that $\mu_2 \gg 1$ and $\mu_3 \ll 1$,
the three-generation constraint is explained \tagP{}/\tagOpen{}.
% BACKFILL-GAP14: μ-window constraint END

\subsection{Mechanism C: Bulk Topology}

If the 5D bulk manifold $\mathcal{M}^5$ has fundamental group $\pi_1(\mathcal{M}^5) = \mathbb{Z}_3$,
winding modes would come in three classes.

\textbf{Verdict:} RED---Pure speculation. EDC specifies local geometry, not global topology.

% ════════════════════════════════════════════════════════════════════════════════
\section{Summary Table}
\label{sec:gen_summary}
% ════════════════════════════════════════════════════════════════════════════════

\begin{center}
\begin{tabular}{lccc}
\toprule
\textbf{Mechanism} & \textbf{Derived?} & \textbf{Explains $n \leq 2$?} & \textbf{Verdict} \\
\midrule
A: $\mathbb{Z}_3 \subset \mathbb{Z}_6$ & Partial \tagI{} & No & YELLOW \\
B: KK truncation & No \tagP{} & No & RED \\
C: Bulk topology & No \tagP{} & No & RED \\
\bottomrule
\end{tabular}
\end{center}

% ════════════════════════════════════════════════════════════════════════════════
\section{Falsifiability}
\label{sec:gen_falsify}
% ════════════════════════════════════════════════════════════════════════════════

\begin{tcolorbox}[edcCanonical, title=\textbf{Falsifiability Statement}]
\textbf{Prediction:} If the $\mathbb{Z}_3$ structure underlies generation count:
\begin{equation}
    N_{\text{gen}} = 3 \quad \text{exactly}
\end{equation}

\textbf{Falsification criterion:}
\begin{itemize}[nosep]
    \item Discovery of a 4th generation fermion $\Longrightarrow$ EDC mechanism fails
    \item A sequential fourth charged lepton $\ell_4^-$ or fourth up-type quark $t'$
          would invalidate the $\mathbb{Z}_3$ identification
\end{itemize}

\textbf{Current experimental status} \tagBL{}:
\begin{itemize}[nosep]
    \item LEP: $N_\nu = 2.984 \pm 0.008$ from $Z$ width
    \item LHC: No evidence for 4th generation quarks up to $\sim 1$ TeV
    \item Precision EW: 4th generation with SM-like couplings strongly disfavored
\end{itemize}

\textbf{Verdict:} Current data \textbf{consistent} with $N_{\text{gen}} = 3$, but
EDC gains predictive power only if the derivation is completed.
\end{tcolorbox}

% ════════════════════════════════════════════════════════════════════════════════
\section{Open Problems}
\label{sec:gen_open}
% ════════════════════════════════════════════════════════════════════════════════

\begin{enumerate}
    \item \textbf{KK truncation calculation (OPR-12):} Compute $V(\xi)$ from EDC action
          and determine mode lifetimes. Show $\tau_2 \gg t_{\text{obs}}$ while
          $\tau_3 \ll t_{\text{obs}}$.

    \item \textbf{$\mathbb{Z}_3$ coupling derivation:} Explain \emph{why} fermion
          wavefunctions couple to the three-fold rotational structure.

    \item \textbf{Bulk topology:} Constrain $\pi_1(\mathcal{M}^5)$ from dynamics,
          or show it's irrelevant for generation structure.

    \item \textbf{Generational mixing:} If generations arise from $\mathbb{Z}_3$,
          what determines CKM/PMNS mixing? (See Ch.~\ref{ch:ckm}, Ch.~\ref{ch:neutrinos}.)
\end{enumerate}

% ════════════════════════════════════════════════════════════════════════════════
\section{Chapter Summary}
\label{sec:gen_chapter_summary}
% ════════════════════════════════════════════════════════════════════════════════

\begin{statusbox}{Key Takeaways}
\begin{enumerate}
\item \textbf{Problem:} SM takes $N_{\text{gen}} = 3$ as input. EDC seeks geometric derivation.

\item \textbf{Identification:} $|\mathbb{Z}_3| = 3$ matches observed generation count \tagI{}.

\item \textbf{Mode structure:} Lepton masses fit tower with $n = 0, 1, 2$ \tagI{}.

\item \textbf{Candidate mechanisms:}
      \begin{itemize}[nosep]
          \item A: $\mathbb{Z}_3$ cardinality---YELLOW (identified)
          \item B: KK truncation---RED (plausible, uncomputed)
          \item C: Bulk topology---RED (speculative)
      \end{itemize}

\item \textbf{Falsifiability:} 4th sequential generation would invalidate $\mathbb{Z}_3$ identification.

\item \textbf{Status:} Generation structure is \textbf{identified} \tagI{} and
      \textbf{postulated} \tagP{}, not \textbf{derived}. This is an open problem.
\end{enumerate}
\end{statusbox}

\begin{tcolorbox}[colback=orange!5, colframe=orange!50!black,
    title=\textbf{Epistemic Audit: Chapter 8}]
\begin{center}
\begin{tabular}{lll}
\toprule
\textbf{Claim} & \textbf{Status} & \textbf{Source} \\
\midrule
$N_{\text{gen}} = 3$ observed & \tagBL{} & PDG, LEP \\
$\mathbb{Z}_6$ from hexagonal packing & \tagP{} & 2D packing (not derived from action) \\
$\mathbb{Z}_6 = \mathbb{Z}_2 \times \mathbb{Z}_3$ & \tagM{} & Group theory \\
$|\mathbb{Z}_3| = 3 \leftrightarrow N_{\text{gen}}$ & \tagI{} & Numerical matching \\
Mode indices $n = 0,1,2$ & \tagI{} & Mass fits \\
KK truncation mechanism & \tagP{} & Not computed \\
No 4th generation prediction & \tagI{}/\tagP{} & Follows from $\mathbb{Z}_3$ \emph{if} derivation holds \\
\bottomrule
\end{tabular}
\end{center}
\end{tcolorbox}

\textbf{Next:} Chapter~\ref{ch:neutrinos} explains neutrino properties as
edge modes with suppressed Higgs overlap.

